\documentclass[12pt,a4paper]{article}
\usepackage[utf8]{inputenc}
\usepackage[T1]{fontenc}
\usepackage{amsmath}
\usepackage[left=2.00cm, right=2.00cm, top=2.00cm, bottom=2.00cm]{geometry}
\usepackage{amssymb}
\usepackage[italian]{babel}

\title{Soluzione Tutorato 1\\Complementi di Matematica per le Scienze Chimiche\\ A.A. 2021/2022}
\author{prof. Emanuele Bottazzi}
\date{}

\usepackage{amsthm}
\theoremstyle{definition}
\newtheorem{exe}{Esercizio}
\newtheorem*{sol}{Soluzione}

\begin{document}
\maketitle

\begin{exe}
    Scrivi (1) una funzione $f: \mathbb{R}^2 \to \mathbb{R}$; (2) una funzione $g: \mathbb{R}^3 \to \mathbb{R}$; (3) una funzione $h: \mathbb{R}^2 \setminus \{(0,0)\} \to \mathbb{R}$.
\end{exe}
\begin{sol}
    Per esempio $f(x,y) = x+y$, $g(x,y,z) = xyz$, $h(x,y) = \frac{x+y}{x^2+y^2}$.
\end{sol}

\begin{exe}
    Rappresenta le curve di livello delle funzioni...
\end{exe}
\begin{sol}
    (1) $f(x,y) = x^2+y^2=c \implies$ circonferenze centrate nell'origine.
    (2) $f(x,y) = 4x^2+y^2=c \implies$ ellissi centrate nell'origine.
    (3) $f(x,y) = x^2-y^2=c \implies$ iperboli.
    (4) $f(x,y) = \sqrt{x^2+y^2}=c \implies$ circonferenze.
    (5) $f(x,y) = \frac{1}{x^2+y^2}=c \implies$ circonferenze.
    (6) $f(x,y) = \frac{y}{x^2-1}=c \implies y = c(x^2-1)$, parabole.
\end{sol}

\begin{exe}
    Rappresenta sul piano cartesiano i punti...
\end{exe}
\begin{sol}
    $(1, \pi) \to (-1, 0)$; $(1,1) \to (\cos 1, \sin 1)$; $(2, 5/4 \pi) \to (-\sqrt{2}, -\sqrt{2})$.
\end{sol}

\begin{exe}
    Calcola il cambiamento di coordinate da cartesiane a polari rispetto al punto $(x_0, y_0)$.
\end{exe}
\begin{sol}
    Il cambiamento di coordinate da cartesiane centrate in $(0,0)$ a polari è $x = \rho \cos \theta, y = \rho \sin \theta$. Invece il cambiamento di coordinate da cartesiane centrate in $(0,0)$ a cartesiane centrate in $(x_0, y_0)$ è $x = x_0 + X, y = y_0 + Y$. Per ottenere il cambiamento di coordinate da cartesiane a polari centrate in $(x_0, y_0)$, passiamo prima a coordinate cartesiane centrate in $(x_0, y_0)$ e poi da queste alle coordinate polari. Otteniamo $x = x_0 + \rho \cos \theta, y = y_0 + \rho \sin \theta$.
\end{sol}

\begin{exe}
    Scrivi funzioni che soddisfano limiti...
\end{exe}
\begin{sol}
    Per esempio $f(x,y) = x+y+2$, $g(x,y,z) = x+y+z-5$.
\end{sol}

\begin{exe}
    Calcola i limiti...
\end{exe}
\begin{sol}
    (1) non esiste; (2) non esiste; (3) $1$; (4) $0$; (5) $0$.
\end{sol}

\begin{exe}[Difficile]
    Spiega perché i seguenti limiti non esistono...
\end{exe}
\begin{sol}
    Perché i limiti lungo diverse direzioni sono diversi.
\end{sol}

\begin{exe}
    Calcola i limiti...
\end{exe}
\begin{sol}
    (1) $1/7$; (2) $+\infty$; (3) non esiste.
\end{sol}

\begin{exe}
    Derivate parziali...
\end{exe}
\begin{sol}
    Per esempio (1) $f(x,y) = \sin(x) \cos(y)$ e $g(x,y) = |y|$.
\end{sol}

\begin{exe}
    Gradiente di $f(x,y) = \frac{xy}{x^2+y^2}$ in $(-1,1)$.
\end{exe}
\begin{sol}
    $\nabla f(-1,1) = (1/2, 1/2)$.
\end{sol}

\begin{exe}
    Derivate parziali di $f(x,y) = \min\{x, y-1\}$.
\end{exe}
\begin{sol}
    $f(0,0)=0 = f(1,1)=0; (0,1); (2,4) \to (0,1); (1,1) \to$ non esiste.
\end{sol}

\end{document}
